% ============================================================
% 开题报告专用格式设置
% 正文主体沿用 nwafuthesis 的章、节、正文、图表与页眉页脚体系。
%
% 封面实现原则：
% 1. Word/PDF 只用于确定字体、字号、页面分区、块间距、缩进与横线长度；
% 2. LaTeX 采用正常段落、居中、minipage、makebox、underline 等结构化排版；
% 3. 不逐字符绝对定位，不使用逐字坐标复刻。
% ============================================================

\graphicspath{{figs/}}

% 学位层次显示文字：由 settings/degree.tex 的 doctor/master 自动决定
\newcommand{\ReportDegreeName}{\ifReportDoctor 博士\else 硕士\fi}

% 占位提示：交稿前建议删除或替换。
\newcommand{\placeholder}[1]{%
  {\color{gray}\small\emph{[待填写：#1]}}%
}

% 图示占位框
\newcommand{\figureplaceholder}[2][0.88\textwidth]{%
  \begin{figure}[htbp]
    \centering
    \fbox{\parbox[c][5.0cm][c]{#1}{\centering\color{gray}#2}}
    \caption{#2}
  \end{figure}
}

% 正文页眉：根据硕/博选择自动切换。
\newcommand{\setreportheader}{%
  \nwafuhead[C]{%
    \small
    \ifodd\value{page}
      \nouppercase{\leftmark}%
    \else
      西北农林科技大学\ReportDegreeName 研究生开题报告%
    \fi
  }%
  \nwafufoot[C]{\small\thepage}%
}

% ============================================================
% 封面字体
% 字体/字号依据原 Word：
%   附件 1：等线 Light 14 pt
%   校名、学位论文、开题报告：宋体 18 pt
%   题目标签：黑体 22 pt
%   信息字段与日期：宋体 15 pt
%   英文/数字：Times New Roman
% ============================================================
% Windows 下优先使用原 Word 对应字体；其他系统使用可稳定编译的开源回退字体。
% 用字体文件是否存在来判断，避免 Linux/fontconfig 将 SimSun 等名称错误映射为替代字体。
\IfFileExists{C:/Windows/Fonts/simsun.ttc}
  {\newCJKfontfamily\CoverSong{SimSun}}
  {\newCJKfontfamily\CoverSong[BoldFont=FandolSong-Bold.otf]{FandolSong-Regular.otf}}
\IfFileExists{C:/Windows/Fonts/simhei.ttf}
  {\newCJKfontfamily\CoverHei{SimHei}}
  {\newCJKfontfamily\CoverHei[BoldFont=FandolHei-Bold.otf]{FandolHei-Regular.otf}}
\IfFileExists{C:/Windows/Fonts/Dengl.ttf}
  {\newCJKfontfamily\CoverDeng{DengXian Light}}
  {\newCJKfontfamily\CoverDeng[BoldFont=FandolHei-Bold.otf]{FandolHei-Regular.otf}}
\IfFileExists{C:/Windows/Fonts/times.ttf}
  {\newfontfamily\CoverTNR{Times New Roman}}
  {\newfontfamily\CoverTNR{lmroman10-regular.otf}}

% ============================================================
% 封面结构化组件
% ============================================================

% 下划线填写框：正常 LaTeX 行内排版，不使用绝对坐标。
\newcommand{\CoverFillLine}[2][100bp]{%
  \underline{\makebox[#1][c]{#2}}%
}

% 中英文题目统一题目块。
% 原 Word 中“题目（居中）：”“英文题目（居中）：”只是空白模板的填写提示。
% 因此：
%   - 对应题目为空时：显示提示文字；
%   - 对应题目填写后：提示文字自动消失，题目占据整行标题区并居中；
%   - 长题目使用 minipage 正常自动换行，每一行保持居中；
%   - 中英文题目宽度统一受 Word 左/右 90 bp 页边距约束；
%   - 不使用 makebox 承载长题目，避免横向溢出。
% 封面页边距按原 Word 的页面设置统一约束：左/右各 1.25 in = 90 bp。
% 题目只允许在这一正文版心宽度内排版，不能进入页边距。
\newlength{\CoverPageMargin}
\newlength{\CoverTitleTextWidth}
\newlength{\CoverChineseTitleBlockWidth}
\newlength{\CoverEnglishTitleBlockWidth}
\newlength{\CoverTitleRegionHeight}
\setlength{\CoverPageMargin}{90bp}
\setlength{\CoverTitleTextWidth}{\dimexpr\paperwidth-2\CoverPageMargin\relax}
\setlength{\CoverChineseTitleBlockWidth}{\CoverTitleTextWidth}
\setlength{\CoverEnglishTitleBlockWidth}{\CoverTitleTextWidth}
% 题目区向上扩展 40bp，用于容纳 Word 中 22pt、1.5 倍行距的长标题。
% “题目区前置间距 + 题目区高度”的总和保持不变，因此基本信息区的位置不会变化。
\setlength{\CoverTitleRegionHeight}{170.6bp}

% 空白模板中的中文题目提示
\newcommand{\CoverChineseTitlePrompt}{%
  \noindent\hspace*{120bp}%
  {\CoverHei\fontsize{22bp}{33bp}\selectfont 题目（居中）：}\par
}

% 空白模板中的英文题目提示
\newcommand{\CoverEnglishTitlePrompt}{%
  \noindent\hspace*{120bp}%
  {\CoverHei\fontsize{22bp}{33bp}\selectfont 英文题目（居中）：}\par
}

% 已填写的中英文题目放在同一个“固定高度题目区”中。
%
% 行距规则：
%   1. 中文题目自身换行：22 bp 字号，行距按填写后的 Word 实际基线间距微调；
%   2. 英文题目自身换行：22 bp 字号、24 bp 固定基线距离，与 Word 填写效果一致；
%   3. 中文题目与英文题目之间保留独立段间距，且不影响后续基本信息位置；
%   4. 题目区总高度仍固定为 \CoverTitleRegionHeight，
%      因此无论题目是一行还是多行，都不会推动后面的基本信息区。
\newlength{\CoverChineseTitleLeading}
\newlength{\CoverEnglishTitleLeading}
\newlength{\CoverBilingualTitleGap}
% 封面标题行距独立于论文正文的全局行距设置。
% 填写后的 Word/PDF 实测：中文两行基线约 31.2bp，英文约 24bp。
\setlength{\CoverChineseTitleLeading}{31.2bp}
\setlength{\CoverEnglishTitleLeading}{24bp}
\setlength{\CoverBilingualTitleGap}{16bp}

\newcommand{\CoverFilledBilingualTitles}{%
  \vbox to \CoverTitleRegionHeight{%
    \vss
    \vspace*{-35bp}% 视觉上移，但不改变固定题目区总高度

    % ---------------- 中文题目 ----------------
    \ifdefempty{\ReportTitle}{}{%
      \noindent\makebox[\paperwidth][c]{%
        \begin{minipage}{\CoverChineseTitleBlockWidth}%
          \centering
          \linespread{1}\selectfont\CoverHei\fontsize{22bp}{\CoverChineseTitleLeading}\selectfont
          \setlength{\parskip}{0pt}%
          \setlength{\emergencystretch}{2em}%
          \sloppy
          \ReportTitle\par
        \end{minipage}%
      }\par
    }%

    % ---------------- 中英文题目之间 ----------------
    % 只有中、英文题目都已填写时才加入该间距。
    % 题间距由独立参数控制。
    \ifboolexpr{
      not test {\ifdefempty{\ReportTitle}}
      and
      not test {\ifdefempty{\ReportTitleEN}}
    }{%
      \vspace*{\CoverBilingualTitleGap}%
    }{}%

    % ---------------- 英文题目 ----------------
    \ifdefempty{\ReportTitleEN}{}{%
      \noindent\makebox[\paperwidth][c]{%
        \begin{minipage}{\CoverEnglishTitleBlockWidth}%
          \centering
          \linespread{1}\selectfont\CoverTNR\fontsize{22bp}{\CoverEnglishTitleLeading}\selectfont
          \setlength{\parskip}{0pt}%
          \setlength{\emergencystretch}{3em}%
          \sloppy
          \ReportTitleEN\par
        \end{minipage}%
      }\par
    }%

    \vss
  }%
}

% 中文 + 英文统一组件，并且总高度固定。
% - 两项都为空：保留原 Word 的“题目（居中）：/英文题目（居中）：”提示位置；
% - 任一题目填写：提示文字消失，题目在固定区域内居中排版；
% - 固定区域结束后直接进入基本信息区，所以基本信息不会被长题目下推。
\newcommand{\CoverBilingualTitleBlock}{%
  \ifboolexpr{ test {\ifdefempty{\ReportTitle}} and test {\ifdefempty{\ReportTitleEN}} }
    {%
      \vbox to \CoverTitleRegionHeight{%
        % 题目区本身向上扩展了 40bp；空白提示仍保持在原来的填写位置。
        \vspace*{40bp}%
        \CoverChineseTitlePrompt
        \vspace*{9.2bp}%
        \CoverEnglishTitlePrompt
        \vfil
      }%
    }%
    {\CoverFilledBilingualTitles}%
}

% ----------------------------------------------------------------
% 基本信息区：统一两列网格
%
% 目标：
%   1. 六行标签使用同一个左边界、同一个右边界；
%   2. 短标签通过 \hfil 在“语义块”之间自动拉伸，形成类似 Word
%      “分散对齐”的效果，而不是手工给每个字塞不同空格；
%   3. 所有填写横线从同一位置开始，并保持完全相同的长度；
%   4. 信息内容仍然是正常 LaTeX 流式排版，不做绝对坐标定位。
% ----------------------------------------------------------------
\newlength{\CoverInfoLeftIndent}
\newlength{\CoverInfoLabelWidth}
\newlength{\CoverInfoGap}
\newlength{\CoverInfoLineWidth}

% 保持横线起点接近原 PDF（约 304.2 bp），同时给最长标签留足宽度。
\setlength{\CoverInfoLeftIndent}{139.7bp}
\setlength{\CoverInfoLabelWidth}{157bp}
\setlength{\CoverInfoGap}{7.5bp}
\setlength{\CoverInfoLineWidth}{172.6bp}

% 标签文字的“分散对齐”：在相邻字符之间加入可伸缩胶，
% 使不同字数的标签都能同时贴齐同一个左边界和右边界。
% 这是正常盒模型排版，不是逐字绝对定位。
\ExplSyntaxOn
\NewDocumentCommand{\CoverDistributedText}{m}
  {
    \bool_set_true:N \l_tmpa_bool
    \tl_map_inline:nn {#1}
      {
        \bool_if:NTF \l_tmpa_bool
          { \bool_set_false:N \l_tmpa_bool }
          { \hfil }
        ##1
      }
  }
\ExplSyntaxOff

% 标签盒：固定宽度 + 字符分散对齐。
\newcommand{\CoverInfoLabelBox}[1]{%
  \hbox to \CoverInfoLabelWidth{{\CoverSong\fontsize{15bp}{18bp}\selectfont
    \CoverDistributedText{#1}}}%
}

% 信息栏专用填写线：所有行强制使用同一长度。
\newcommand{\CoverInfoValueLine}[1]{%
  {\CoverSong\fontsize{15bp}{18bp}\selectfont
   \CoverFillLine[\CoverInfoLineWidth]{#1}}%
}

% 统一信息行：标签列 + 固定间距 + 等长填写线。
\newcommand{\CoverInfoRow}[2]{%
  \noindent\hspace*{\CoverInfoLeftIndent}%
  \CoverInfoLabelBox{#1}%
  \hspace*{\CoverInfoGap}%
  \CoverInfoValueLine{#2}%
}

% ----------------------------------------------------------------
% 开题论证小组成员：由用户手工决定换行
%
% 使用方式：
%   \newcommand{\ReportGroupMembers}{%
%     张三 李四 王五\\
%     赵六 孙七 周八%
%   }
%
% 规则：
%   1. 用 \\ 明确指定在哪里换到下一条横线；
%   2. 每一行都在等长横线上居中；
%   3. 至少保留两条成员填写线（与原 Word 一致）；
%   4. 如需第 3、4 行，继续写 \\ 即可；
%   5. 模板不再依据姓名长度自动改变你的分行决定。
% ----------------------------------------------------------------
\ExplSyntaxOn
% 同一行不同姓名之间的固定间距。
% 用户只需要输入普通空格；模板会把一个或多个空格归一为一个不可断开的标准空格。
\newcommand{\MemberSep}{\nobreakspace}

\seq_new:N \l__nwafu_member_manual_lines_seq
\seq_new:N \l__nwafu_member_manual_raw_lines_seq
\tl_new:N  \l__nwafu_member_manual_source_tl
\tl_new:N  \l__nwafu_member_manual_line_tl
\int_new:N \l__nwafu_member_manual_line_count_int

% 将用户填写内容完整展开后：
% 1. 只以显式的 \\ 为换行分隔符；
% 2. 每一行内部，一个或多个普通空格自动替换为 \MemberSep；
% 3. 行首、行尾空格自动去除。
\cs_new_protected:Npn \__nwafu_member_manual_build:n #1
  {
    \seq_clear:N \l__nwafu_member_manual_lines_seq
    \seq_clear:N \l__nwafu_member_manual_raw_lines_seq
    \tl_set:Nx \l__nwafu_member_manual_source_tl {#1}
    \seq_set_split:NnV \l__nwafu_member_manual_raw_lines_seq { \\ } \l__nwafu_member_manual_source_tl

    \seq_map_inline:Nn \l__nwafu_member_manual_raw_lines_seq
      {
        \tl_set:Nn \l__nwafu_member_manual_line_tl {##1}
        \tl_trim_spaces:N \l__nwafu_member_manual_line_tl
        % 连续空格也只作为一个“姓名分隔符”，显示为一个标准空格。
        \regex_replace_all:nnN { \s+ } { \c{MemberSep} } \l__nwafu_member_manual_line_tl
        \seq_put_right:NV \l__nwafu_member_manual_lines_seq \l__nwafu_member_manual_line_tl
      }

    % 原 Word 默认预留两行；不足两行时自动补空行。
    \int_set:Nn \l__nwafu_member_manual_line_count_int
      { \seq_count:N \l__nwafu_member_manual_lines_seq }
    \int_while_do:nn { \l__nwafu_member_manual_line_count_int < 2 }
      {
        \seq_put_right:Nn \l__nwafu_member_manual_lines_seq { }
        \int_incr:N \l__nwafu_member_manual_line_count_int
      }
  }

% 输出每一条“等长横线 + 居中姓名”。
% 日期的位置是固定的，因此成员行数较多时，只压缩成员行之间的垂直间距，
% 不再把日期区向下推。手工 \ 仍然只决定“在哪里换行”。
\dim_new:N \l__nwafu_member_manual_gap_dim
\cs_new_protected:Npn \__nwafu_member_manual_typeset:
  {
    \int_set:Nn \l__nwafu_member_manual_line_count_int
      { \seq_count:N \l__nwafu_member_manual_lines_seq }
    \int_case:nnF { \l__nwafu_member_manual_line_count_int }
      {
        {1}{\dim_set:Nn \l__nwafu_member_manual_gap_dim {13.8bp}}
        {2}{\dim_set:Nn \l__nwafu_member_manual_gap_dim {13.8bp}}
        {3}{\dim_set:Nn \l__nwafu_member_manual_gap_dim {8.0bp}}
        {4}{\dim_set:Nn \l__nwafu_member_manual_gap_dim {3.0bp}}
      }
      {\dim_set:Nn \l__nwafu_member_manual_gap_dim {1.0bp}}
    \seq_map_indexed_inline:Nn \l__nwafu_member_manual_lines_seq
      {
        \CoverInfoValueLine{##2}\par
        \int_compare:nNnF {##1} = { \seq_count:N \l__nwafu_member_manual_lines_seq }
          { \vspace*{\l__nwafu_member_manual_gap_dim} }
      }
  }

\NewDocumentCommand{\CoverGroupMembersBlock}{m}
  {
    \__nwafu_member_manual_build:n {#1}
    \noindent\hspace*{\CoverInfoLeftIndent}%
    \begin{minipage}[t]{\CoverInfoLabelWidth}
      \CoverInfoLabelBox{开题论证小组成员姓名}%
    \end{minipage}%
    \hspace*{\CoverInfoGap}%
    \begin{minipage}[t]{\CoverInfoLineWidth}
      \__nwafu_member_manual_typeset:
    \end{minipage}%
  }
\ExplSyntaxOff

% 日期行内容。日期块的页面位置在 \makeopeningcover 中固定，
% 因而不会受题目行数或成员行数影响。
\newcommand{\CoverDateContent}{%
  {\CoverSong\fontsize{15bp}{18bp}\selectfont
    研究生通过开题论证日期：%
    \hspace{5.4bp}\CoverFillLine[20bp]{{\CoverTNR\ReportDefenseYear}}年%
    \hspace{5.4bp}\CoverFillLine[20bp]{{\CoverTNR\ReportDefenseMonth}}月%
    \hspace{5.4bp}\CoverFillLine[20bp]{{\CoverTNR\ReportDefenseDay}}日%
  }%
}

% ============================================================
% 一页式封面
% 只按“排版块”复现：整体缩进、行高、块间距和横线宽度来自 PDF 测量，
% 文字本身仍由 LaTeX 正常组版，不进行逐字符定位。
% ============================================================
\newcommand{\makeopeningcover}{%
  \clearpage
  % 封面临时使用整页版心。geometry 的切换必须发生在局部分组之外：
  % 这样封面中的普通 LaTeX 排版才真正以纸张左边界为基准，
  % 同时也能在封面结束后可靠恢复论文模板的目录/正文版心。
  \newgeometry{left=0pt,right=0pt,top=0pt,bottom=0pt}%
  \begingroup
  \thispagestyle{empty}%
  \setlength{\parindent}{0pt}%
  \setlength{\parskip}{0pt}%
  \xeCJKsetup{PunctStyle=plain}%

  % 附件 1
  \vspace*{61.4bp}%
  \noindent\hspace*{90bp}%
  {\CoverDeng\bfseries\fontsize{14bp}{14bp}\selectfont 附件 1}\par

  % 顶部标题区
  \vspace*{43.7bp}%
  \noindent\makebox[\paperwidth][c]{{\CoverSong\fontsize{18bp}{18bp}\selectfont 西北农林科技大学}}\par
  \vspace*{11.3bp}%
  \noindent\makebox[\paperwidth][c]{{\CoverSong\fontsize{18bp}{18bp}\selectfont
    \CoverFillLine[54bp]{{\CoverTNR\ReportGrade}}级\ReportDegreeName 研究生学位论文}}\par
  \vspace*{11.3bp}%
  \noindent\makebox[\paperwidth][c]{{\CoverSong\fontsize{18bp}{18bp}\selectfont 开题报告}}\par

  % 题目区
  \vspace*{62.3bp}%
  \CoverBilingualTitleBlock
  \vspace*{11.0bp}%

  % 基本信息区：题目区总高度已经固定，直接进入信息区。
  \CoverInfoRow{学院（系、所）}{\ReportDepartment}\par
  \vspace*{11.3bp}%
  \CoverInfoRow{一级学科/类别（领域）}{\ReportDiscipline}\par
  \vspace*{11.3bp}%
  \CoverInfoRow{研究生姓名及学号}{\ReportStudent\quad{\CoverTNR\ReportStudentID}}\par
  \vspace*{11.3bp}%
  \CoverInfoRow{指导教师姓名}{\ReportSupervisor}\par
  \vspace*{11.3bp}%
  \CoverInfoRow{开题论证小组组长姓名}{\ReportGroupLeader}\par
  \vspace*{11.3bp}%
  % 成员姓名只按用户写入的 \ 手工换行；普通空格仅分隔姓名；至少保留两条等长填写线。
  \CoverGroupMembersBlock{\ReportGroupMembers}\par

  % 日期区采用“区块级固定锚点”：只固定整行日期的位置，
  % 不是逐字坐标。这样长题目、多人名单都不会改变日期位置。
  \begin{tikzpicture}[remember picture,overlay]
    \node[anchor=north west,inner sep=0pt]
      at ([xshift=146.7bp,yshift=-738.0bp]current page.north west)
      {\CoverDateContent};
  \end{tikzpicture}%

  % 注意：\restoregeometry 不能放在上面的局部分组中。
  % 否则分组结束后，geometry 的页面参数可能重新退回封面的 0 边距，
  % 从而使目录/正文的可用高度异常增大，造成目录项越过页底。
  \endgroup
  \clearpage
  % 在局部分组之外恢复 nwafuthesis 类预先设置的论文版心。
  % doctor/master 对应的版心参数由类文件统一管理，目录与正文不再另行覆盖。
  \restoregeometry
}

\newcommand{\makeopeningapproval}{}
